% !TEX root = ../CerdaMS.tex

\section{A resource calculus with memory}
\label{sec:resource}

%\comrci{With memory or with memories?}

In this section we introduce the \emph{\lamI-resource calculus} refining the (finite) resource calculus~\cite{EhrhardR03}, best known as the target language of Ehrhard and Regnier's Taylor expansion~\cite{Ehrhard.Reg.08}. 
So, the resource calculus is not meant to be a stand-alone language, but rather another theory of approximations for the \lam-calculus.
Before going further, we recall its main properties.

%\subsection{An introduction to resource calculi}

We consider here the promotion-free fragment of the resource calculus introduced in \cite{PaganiT09}.
Its syntax is similar to the \lam-calculus, except for the applications that are of shape $\rt{s}\rb t$, where $\rb t = [\rt{t}_1,\dots,\rt{t}_n]$ is a multiset of resources called \emph{bag}.
The resources populating $\rb t$ are linear as they cannot be erased or copied by $\rt s$, they must be used \emph{exactly once} along the reduction. 
When contracting a term of the form $s = (\lam x.\rt{s'})[\rt{t}_1,\dots,\rt{t}_n]$ there are two possibilities.
\begin{enumerate}
\item If the number of occurrences of $x$ in $\rt{s'}$ is exactly $n$, then each occurrence is substituted by a different $\rt{t}_i$. 
Since the elements in the bag are unordered, there is no canonical bijection between the resources and the occurrences of $x$. 
The solution consists in collecting all possibilities in a formal sum of terms, the sum representing an inner-choice operator.
\item If there is a mismatch between the number of occurrences and the amount of resources, then $s$ reduces to the empty-sum, $\esum{}$. From a programming-language perspective, this can be thought of as a program terminating abruptly after throwing an uncaught exception.
\end{enumerate}
The first-class citizens of the resource calculus are therefore \emph{finite sums} of resource terms, that are needed to ensure the (strong) confluence of reductions. Another important property is strong normalization, that follows from the fact that no resource can be duplicated.

\subsection{\lamI-resource expressions and \lamI-resource sums}

Our version of the resource calculus is extended with labels representing the memory of free variables that were present in the \lamI-term they approximate.
Just like in \cref{def:approximants}\eqref{def:approximants1} 
we endowed the constant $\bot$ with a finite set $X$ of variables, 
here we annotate:
\begin{itemize}
\item the empty bag $1$ of resource terms, since an empty bag of approximants of $M$ should remember the free variables of $M$;
\item the empty sum $\esum{}$ of resource terms. Indeed, if a resource term vanishes during reduction because of the mismatch described above, its free variables should be remembered.
\end{itemize}
%We now rigorously define the terms of our modified resource calculus, 
%as well as the bags used as arguments of applications.
%Then, we extend this syntax to finite sums of terms, 
%on which the dynamics of the calculus is acting.

\begin{definition} \label{def:rterms}\mbox{}
\bsub\item
	For all $X \subf \Var$, the sets $\rterms(X)$ and $\rterms*(X)$
	are defined by induction:
	\begin{gather*}
		\infer{ x \in \rterms(\set{x}) }{} \qquad
		\infer{ \lam x.\rt{s} \in \rterms(X \setminus \set{x}) }
			{ \rt{s} \in \rterms(X) & x \in X } \qquad
		\infer{ \rt{s}\rb{t} \in \rterms(X \cup Y) }
			{ \rt{s} \in \rterms(X) & \rb{t} \in \rterms*(Y) } \\
		\infer{ 1_X \in \rterms*(X) }{} \qquad
		\infer{ [\rt{t}_0,\dots,\rt{t}_n] \in \rterms*(X)}
			{ \rt{t}_0 \in \rterms(X) & \cdots & \rt{t}_n \in \rterms(X) }
	\end{gather*}
\item
	The set $\rterms$ of \emph{\lamI-resource terms} and
	the set $\rterms*$ of \emph{bags} are given by:
	\[	\rterms \eqdef \bigdunion_{X \subf \Var} \rterms(X)
		\quad \text{and}\quad
		\rterms* \eqdef \bigdunion_{X \subf \Var} \rterms*(X). \]
\esub
\end{definition}
As a matter of notation, we let $\rterms**$ denote either $\rterms$ or $\rterms*$, indistinctly but coherently.
We call \emph{resource expressions} generic elements $s,t\in\rterms**$. 
We denote the union of two bags $\rb{t}, \rb{u}\in\rterms*(X)$ multiplicatively by $\rb{t} \cdot \rb{u}$,
whose neutral element is the empty bag, $1_X$.
%whence the empty bag $1_X$ is the neutral element.

\Needspace{2cm}
\begin{remark}\label{rem:aboutresourceterms}\mbox{}
\bsub
\item
In every \lamI-resource term of the form $\lam x.s$, the variable $x$ must occur freely in $s$: it may appear in the undecorated underlying term, or in the \enquote{memory} $X$ decorating $1_X$.
Therefore, it makes sense to define $\fv(\rt s) \eqdef X$ whenever $\rt{s} \in \rterms**(X)$.
\item\label{rem:aboutresourceterms2}
Each set $\rterms*(X)$ is isomorphic to the monoid of
multisets of elements of $\rterms(X)$.
Notice that $\rterms*$ is \emph{not} the set of all bags
of elements of $\rterms$, just its subset of bags
whose elements \emph{have the same free variables}
(and so inductively in the subterms).
\esub
\end{remark}

\begin{example}\label{ex:resourceterms}\mbox{}
\bsub
\item The identity $\comb{I}$ belongs to $\rterms$, whereas the projections do not: $\comb{K},\comb{F}\notin\rterms$.
\item Note that $\lam xy.x1_{\emptyset}\notin\rterms$ since $y\notin\fv(x1_{\emptyset})$, but $\lam xy.x1_{\set{y}}\in\rterms$ because $y\in\set{y}$.
\item The terms $\comb{D}_0 = \lam x.x1_{\set{x}}$ and $\comb{D}_{n+1} = \lam x.x[x,\ldots,x]$, where the bag contains $n+1$ occurrences of $x$, are \lamI-resource terms. By Remark~\ref{rem:aboutresourceterms}\eqref{rem:aboutresourceterms2}, we obtain:
\item $[x,y] \notin \rterms*$, while $[x, x[x], x[\lam y.y,\lam y.y[y]], (\lam y.y)[x],  \lam y.y1_{\set x}] \in \rterms*(\set{x})\subseteq\rterms*$.
\esub
\end{example}

We consider resource expressions up to $\alpha$-equivalence, under the proviso that abstractions bind linear occurrences of variables as well as occurrences in the memory of empty bags. For instance, $\lam xyz.x1_{\set{x}}1_{\set{x,y,z}}$ and  $\lam x'y'z.x'1_{\set{x'}}1_{\set{x',y',z}}$ are considered $\alpha$-equivalent. 

\begin{definition}\label{def:frsums}
	For all $X \subf \Var$, the set $\frsums[X]**$
	of \emph{sums of \lamI-resource terms} (\enquote{\mbox{resource sums}}, for short)
	is defined as the $\Nat$-semimodule
	of finitely supported formal sums of expressions in $\rterms**(X)$,
	with coefficients in $\Nat$.
	Explicitly, it can be presented as:
	\bnf[ \frsums[X]** ]{ \rs{s},\rs{t} }
		{ \esum X \mid \rt r\mid \rt{r} + \rs{s} }
		[ \text{for $\rt{r} \in \rterms**(X)$} ]
	quotiented by associativity and commutativity of $+$,
	as well as neutrality of $\esum X$.
\end{definition}

Note that $\rterms**(X)\subseteq\frsums[X]**$, that is, resource expressions are assimilated to the corresponding
one-element sum.
The constructors of the calculus are extended
to resource sums by (bi)linearity, \emph{i.e.}
for $\rt{s} \in \rterms(X)$, $\rs{s} \in \frsums[X]$, 
$\rb{t} \in \rterms*(Y)$, $\rbs{t} \in \frsums[Y]*$,
$\rt{u} \in \rterms(Y)$ and $\rs{u} \in \frsums[Y]$, we have:
\[
	\begin{array}{llll}
	(\rt{s}+\rs{s}) \rb{t} \eqdef \rt{s}\rb{t} + \rs{s}\rb{t},&
	\esum X \rb{t} \eqdef \esum {X \cup Y},&
	\lam x.(\rt{s}+\rs{s}) \eqdef \lam x.\rt{s} + \lam x.\rs{s},&
	\lam x.\esum X \eqdef \esum {X \setminus \set{x}},\\[1ex]
	\rt{s} (\rb{t}+\rbs{t}) \eqdef \rt{s}\rb{t} + \rt{s}\rbs{t},&
	\rt{s}\esum Y \eqdef \esum {X \cup Y},&
	[\rt{u} + \rs{u}] \cdot \rb{t} \eqdef [\rt{u}]\cdot\rb{t}
		+ [\rs{u}]\cdot\rb{t},&	
	[\esum Y] \cdot \rb{t} \eqdef \esum Y.
	\end{array}
\]
Therefore, if $\esum X$ occurs in $s$ not as a summand but as a proper subterm, then $s = \esum{\fv(s)}$.

\subsection{Memory substitution and resource substitution}

While the usual finite resource calculus is completely linear, the variables we store in the memory of $1_X$ and $\esum X$ are not. 
The memory $X$ remembers the variables present in $X$ not their amounts, this is the reason why it is modelled as a set, not as a multiset. 
This consideration leads us to define two kinds of substitutions.

The (non-linear) \emph{memory substitution} of a set $Y\subf\Var$ for a variable $x$ in a \lamI-resource term $s$ does not interact with the linear occurrences of $x$ in $s$ (\emph{i.e.}, they remain unchanged), it just replaces the \enquote{memory} of $x$ in the empty bags with the content of $Y$.

\begin{definition} \label{def:msubst}\mbox{}
\bsub
\item For all $X,Y \subf \Var$ and $x \in \Var$, define
	\[
	X \msubst Y \eqdef \begin{cases}
		X \setminus \{x\} \cup Y, 
			& \text{if $x \in X$,} \\
		X, & \text{otherwise.}\\
	\end{cases}
	\]
	\item Given $\rt{s} \in \rterms**(X)$, the \emph{memory substitution} of  $x$ by $Y\subf\Var$ in $\rt{s}$ is the resource term $\rt{s} \msubst Y$ defined as follows (for $x \neq y$ and, in the abstraction case, $y \notin Y$):
	\begin{flalign*}
		x \msubst Y & \eqdef x, &
		1_X \msubst Y & \eqdef 1_{X\msubst Y},&
		(\rt{s}\rb{t}) \msubst Y & \eqdef 
			(\rt{s}\msubst Y)(\rb{t}\msubst Y),  \\
		y  \msubst Y & \eqdef y,  &
		(\lam y.\rt{s}) \msubst Y & \eqdef \lam y.\rt{s} \msubst Y,	&
		([\rt s] \cdot \rb t) \msubst Y & \eqdef
			[\rt s \msubst Y] \cdot (\rb t \msubst Y).
	\end{flalign*}
\esub			
\end{definition}

We now define the \emph{resource substitution} of a bag $\rb u$ for $x$ in $\rt s$, whose effect is twofold.
1)~It non-deterministically replaces each linear occurrence of $x$ with a 
resource from the bag (as usual).
2) It applies the memory substitution of $x$ by $\fv(\rb u)$ to the resulting sum of terms.

\begin{definition} \label{def:rsubst}
	Given $\rt s \in \rterms**(X)$, $x \in \Var$ and $\rb u \in \rterms*(Y)$,
	the \emph{resource substitution} of $x$ by $\rb u$ in $\rt s$
	is the resource sum $\rt s \rsubst u \in \frsums[X \msubst Y]**$
	defined as follows:
	\begin{align*}
		x \rsubst u & \eqdef 
			\begin{cases}
			\rt u,& \text{if $\rb u = [\rt u]$},\\
			\esum Y, & \text{otherwise},
			\end{cases} &
		(\rt{s}\rb{t}) \rsubst u & \eqdef 
			\sum_{\rb u = \rb v \cdot \rb w}
			(\rt{s}\rsubst v)( \rb{t}\rsubst w),\\
		y  \rsubst u & \eqdef 
			\begin{cases}
			y, & \text{if $\rb u = 1_Y$,} \\
			\esum {\set y}, & \text{otherwise},
			\end{cases}&
		1_X \rsubst u & \eqdef 
			\begin{cases}
			1_{X \msubst Y}, & \text{if $\rb u = 1_Y$}, \\
			\esum {X \msubst Y}, & \text{otherwise,} \\
			\end{cases}\\
		(\lam y.\rt{s}) \rsubst u & \eqdef 
			\lam y.\rt{s} \rsubst u, &
		([\rt s] \cdot \rb t) \rsubst u & \eqdef
			\sum_{\rb u = \rb v \cdot \rb w}
			[\rt s \rsubst v] \cdot (\rb t \rsubst w).
	\end{align*}
	with $x \neq y$ and in the abstraction case $y \notin Y$.
	We extend it to sums in $\frsums[X]**$ by setting:
	\begin{align*}
		\esum X \rsubst u & \eqdef \esum {X\msubst Y} &
		(\rt{s} + \rs{s}) \rsubst u & \eqdef 
			\rt{s} \rsubst u + \rs{s} \rsubst u \\
		\rt s \rsubst* {\esum Y} & \eqdef \esum {X \msubst Y} &
		\rt s \rsubst* {(\rb u + \rbs u)} & \eqdef
			\rt s \rsubst u + \rt s \rsubst* {\rbs u}.
	\end{align*}
\end{definition}

It is easy to verify that the definition above does indeed define a resource sum in $\frsums[X \msubst Y]**$, and that it is stable under the quotients of Definition~\ref{def:frsums}.

\begin{definition}The \emph{linear degree} $\deg_x(\rt s)$ of $x \in \Var$
	in some $\rt s \in \rterms**$ is defined by:
	\begin{align*}
	\deg_x(x) & \coloneqq 1, &
	\deg_x(\lam y.\rt s) & \coloneqq \deg_x(\rt s)
		\text{, wlog. $x\neq y$,} &
	\deg_x(1_X) & \coloneqq 0, \\
	\deg_x(y) & \coloneqq 0 ,&
	\deg_x(\rt s \rb t) & \coloneqq \deg_x(\rt s) + \deg_x(\rb t), &
	\deg_x([\rt s] \cdot \rb t) & \coloneqq \deg_x(\rt s) + \deg_x(\rb t).
	\end{align*}

\end{definition}

The next lemma is not strictly needed, but helps understanding resource substitution.

\begin{lemma}\label{lem:aboutressubst}
	Consider $\rt s \in \rterms**(X)$, $x \in \Var$ and
	$\rb u \in \rterms*(Y)$,
	and write $n \coloneqq \deg_x(\rt s)$. Then
	\[	\rt s \rsubst u = 
		\begin{cases}
		\sum\limits_{\sigma \in \mathfrak{S}(n)}
			\rt s\left[ \rt u_{\sigma(1)}/x^{(1)}, \dots, 
			\rt u_{\sigma(n)}/x^{(n)} \right] \msubst Y,
			& \text{if the cardinality of $\rb u$ is $n$,} \\
		\esum {X \msubst Y}, & \text{otherwise,}\\
		\end{cases} 
	\]
	where: $\mathfrak{S}(n)$ is the set of permutations of $\set{1,\dots,n}$;
	$\rt u_1,\dots,\rt u_n$ is any enumeration of the elements in $\rb u$;
	$x^{(1)},\dots,x^{(n)}$ enumerate the occurrences of $x$ in $\rt s$; and
	$\rt s\left[ \rt u_{\sigma(1)}/x^{(1)}, \dots, \rt u_{\sigma(n)}/x^{(n)} \right]$
	is the \lamI-resource term obtained by substituting each element $\rt u_{\sigma(i)}$ for the occurrence $x^{(i)}$.
\end{lemma}
	
	\begin{proof}
	Straightforward induction on $\rt s$.
	\end{proof}

The next Substitution Lemma concerns the commutation of resource substitutions, and is the analogous of \cite[Lemma~2.1.16]{Bare}.
Notice that the assumptions on the substituted variables
are stronger than in the usual resource calculus
\cite[Lemma~2]{Ehrhard.Reg.08},
where only $x \notin Y$ is required.

\newcounter{lem:resource-substitution}
\setcounter{lem:resource-substitution}{\thetheorem}
\begin{lemma}
\label{lem:resource-substitution}
	Given $\rt s \in \rterms**(X)$,
	$\rb u \in \rterms*(Y)$,
	$\rb v \in \rterms*(Z)$,
	and
	$x \in X \setminus Z$, $y \in X \cup Y$,
	\[	\rt s \rsubst u \rsubst v [y] =
		\sum_{\rb v = \rb v' \cdot \rb v''} \rt s \rsubst* {\rb v'} [y]
		\rsubst*{ \rb u \rsubst* {\rb v''}[y] }. \]
\end{lemma}
\begin{proof} By structural induction on $s$, using Lemma~\ref{lem:aboutressubst} in order to apply the induction hypothesis only to those subterms where the hypotheses are actually satisfied.
\end{proof}

\subsection{The operational semantics}

We finally endow resource expressions and resource sums with a notion of reduction $\rred$.

\begin{definition} \label{def:rred}\mbox{}
\bsub
\item
	For each $X \subf \Var$, define the \emph{resource reduction} as a relation between \lamI-resource terms and resource sums, i.e.\
	$\mathord{\rred} \subseteq \rterms**(X) \times \frsums[X]**$:\\[1ex]
	$\infer{ (\lam x.\rt s)\rb t \rred \rt s \rsubst t }{} \qquad
		\infer{ \lam x.\rt s \rred \lam x.\rs s' }{ \rt s \rred \rs s' } \qquad
		\infer{ \rt s \rb t \rred \rs s' \rb t }{ \rt s \rred \rs s' } \qquad
		\infer{ \rt s \rb t \rred \rt s \rbs t' }
			{ \rb t \rred \rbs t' } \qquad
%		\infer{ \rt s \cdot \rb t \rred \rs s' \cdot \rb t }
		\infer{ [\rt s] \cdot \rb t \rred {[\rs s']} \cdot \rb t }
			{ \rt s \rred \rs s' } $\smallskip
\item\label{def:rred3} We extend the reduction relation $\rred$ to resource sums $\frsums[X]** \times \frsums[X]**$, and simultaneously introduce its reflexive closure $\rredr$, as follows: 
%by saying that $\rt s \rred \rs s' $ and $ \rs t \rredr \rs t'$ entail $\rt s + \rs t \rred \rs s' + \rs t'$.
\[
	\infer{\rt s + \rs t \rred \rs s' + \rs t' }{
		\rt s \rred \rs s'
		&
		\rs t \rredr \rs t'
		}
	\qquad
	\infer{\rs t \rredr \rs t'}{\rs t \rred \rs t'}	
	\qquad
	\infer{\rs t \rredr \rs t}{}
\]
%\item
%	We denote by $\rredr$ the reflexive closure of $\rred$. % and by $\redd[\resource]$ its reflexive-transitive closure. 
%\item\label{def:rred3}
%	We extend the reduction relation $\rred$ to resource sums $\frsums[X]** \times \frsums[X]**$ by saying that $\rt s \rred \rs s' $ and $ \rs t \rredr \rs t'$ entail $\rt s + \rs t \rred \rs s' + \rs t'$.
\item
	We denote by $\redd[\resource]$ the reflexive and transitive closure of the relation $\rred$ given in \eqref{def:rred3}. 
\esub
\end{definition}

Observe that $\rred$ is well-defined because $\fv((\lam x.\rt s)\rb t) = \fv(\rt s \rsubst t)$.
Indeed $\fv((\lam x.\rt s)\rb t) = 
\fv(\rt s) \setminus \set x \cup \fv(\rb t)$
with $x \in \fv(\rt s)$, whence $\fv(\rt s \rsubst t) = \fv(\rt s) \msubst {\fv(\rb t)}
= \fv(\rt s) \setminus \set x \cup \fv(\rb t)$.

\begin{example} We use the resource \lamI-terms $\comb{I}$ and $\comb{D}_n$ from Example~\ref{ex:resourceterms}.
\bsub
\item $\comb{D}_0[z]\rred (x1_{\set{x}})\langle[z]/x\rangle = 
	(x\langle[z]/x\rangle)(1_{\set{z}}) + (x\langle1_{\set{z}}/x\rangle)(1_{\set x}\langle [z]/x\rangle) = z1_{\set z}+\esum{\set{z}} = z1_{\set z}$. 
	Similarly, $\comb{D}_0[\comb{I}] \rred \comb{I}1_{\emptyset} 
	\red \esum{\emptyset}$.
%	, while $\comb{D}_1[\comb{I},\comb{I}]\rred \comb{I}[\comb{I}] 
%	\rred \comb{I}$ is a non-zero reduction. 
\item $(\lam x.\comb{D_0}[x])[z]$ has two redexes. Contracting the outermost first gives $(\lam x.\comb{D_0}[x])[z]\rred \comb{D_0}[z]\rred z1_{\set z}$.
Contracting the innermost $(\lam x.\comb{D_0}[x])[z]\rred (\lam x.x1_{\set x})[z]\rred z1_{\set z}$.
\item $\comb{D}_1[\comb{I},\comb{I}]\rred \comb{I}[\comb{I}] + \comb{I}[\comb{I}]\rred \comb{I}+ \comb{I}[\comb{I}]\rred \comb{I}+\comb{I} = 2.\comb{I}$. 
Thus, sums can arise from single terms.
\esub
\end{example}

We show that $\rred$ enjoys the properties of strong confluence and strong normalisation.

\newcounter{the:rred-norm-confl:confl}
\setcounter{the:rred-norm-confl:confl}{\thetheorem}
\begin{theorem} \label{the:rred-norm-confl}\mbox{}
\bsub
\item The reduction $\rred$ is strongly normalising.
	\label{the:rred-norm-confl:norm}
\item $\rred$ is strongly confluent in the following sense:
	for all $\rs s, \rs t_1, \rs t_2 \in \frsums[X]**$
	such that $\rs s \rred \rs t_1$ and $\rs s \rred \rs t_2$,
	there exists $\rs u \in \frsums[X]**$ such that
	$\rs t_1 \rredr \rs u$ and $\rs t_2 \rredr \rs u$.
	\label{the:rred-norm-confl:confl}
\esub 
\end{theorem}

	\begin{proof}\mbox{}
	\bsub
	\item 
	The size $\size s\in\Nat$ of a resource expression $s\in\rterms**$ 
	is defined by structural induction as usual, 
	with base cases $\size{x} = 1$ and $\size{1_X} = 0$, 
	so that $\size{1_X\cdot \rb t} = \size{\rb t}$.
	The \emph{sum-size} of a resource sum $\rs s$ is 
	the finite multiset defined by $\sumsize{\esum{X}} = []$ and 
	$\sumsize{s+\rs t} = [\size{s}]\cdot \sumsize{\rs t}$.
	Sum-sizes are well-ordered by the usual well-founded ordering 
	$\msetlt$ on finite multisets over $\Nat$ 
	(see \cite[\S A.6]{Terese}).
	
	Now, assume that $\rs s \rred\rs s'$. 
	Since the contraction of a redex suppresses an abstraction 
	and there is no duplication, 
	we get $\sumsize{\rs s}\msetlt \sumsize{\rs s'}$. 
	Conclude since $\msetlt$ is well-founded.
	\item The sequence of lemmas leading to the result follows
	a well-trodden path in the resource calculus \cite{Vaux.19,Cerda.24}.
	It is sufficient to check that they generalize to our setting.
%	We give some details in the appendix.
	\qedhere\esub
	\end{proof}

By \eqref{the:rred-norm-confl:confl} above every $\rs s \in \frsums[X]**$ has a unique $\mathrm{r}$-normal form which is denoted $\nf(\rs s)$. 
%We use the notation $t\in\nf(\rs s)$ to express the fact that $\nf(\rs s) = t + \rs s'$, for some $\rs s'$.


